\documentclass[11pt]{article}
\usepackage[margin=0.92in]{geometry}
\usepackage{amsmath,amssymb,amsthm,mathtools}
\usepackage[hidelinks]{hyperref}

\newtheorem{theorem}{Theorem}
\newtheorem{proposition}{Proposition}
\newtheorem{remark}{Remark}
\newcommand{\D}{\mathbb D}
\newcommand{\C}{\mathbb C}
\newcommand{\N}{\mathbb N}

\title{A Counterexample to Formal Elementary-Symmetric Convergence\\
in the Analytic Wiener Algebra}
\author{Run fa\_banach\_001, agent\_lane\_16}
\date{2026-08-12}

\begin{document}
\maketitle

\section*{Source question and answer}

For the elementary symmetrization map
\[
 \pi:\D^d\longrightarrow\C^d,
 \qquad \pi(z)=(s_1(z),\ldots,s_d(z)),
\]
Section 7.1 of Sasane, arXiv:2201.00183, asks the following.  Given
$f\in W^+_{\rm sym}(\D^d)$, decompose $f=\sum_{k\ge0}p_k$ into
homogeneous symmetric polynomials and write each $p_k$ as a polynomial
in $s_1,\ldots,s_d$.  Does the resulting formal series
$g\in\C[[s_1,\ldots,s_d]]$ converge on a domain
containing $\pi(\D^d)$, so that $f=g\circ\pi$?

Under the standard meaning of convergence of a multivariable power series,
the answer is \emph{no}.  In fact, the illustrative function printed
immediately after the question is already a counterexample.  There is still
a corrected positive result: retaining the transformed polynomials as
weighted-homogeneous blocks always gives a locally uniformly convergent
holomorphic factor on $\pi(\D^d)$.

\section*{Proof Intuition}

The original Wiener norm sees the coefficient mass of a symmetric polynomial
\emph{before} changing coordinates.  Rewriting in $s_1,s_2$ can create large
coefficients which cancel inside each polynomial block.  Along the diagonal
$z=w=r$, the block values remain small, but taking absolute values of their
individual monomials replaces cancellation by a binomial sum.  That sum has
exponential base $3r^2$, which exceeds one well inside the bidisc.  A power
series cannot converge on an open set through such a point.

\begin{theorem}[full negative answer in dimension two]\label{thm:counter}
Let
\begin{equation}\label{eq:f}
 f(z,w)=\sum_{n=1}^{\infty}\frac{(z^2+w^2)^n}{n^2 2^n}.
\end{equation}
Then $f\in W^+_{\rm sym}(\D^2)$.  Under $s_1=z+w$, $s_2=zw$, the
formal series prescribed by the source is
\begin{equation}\label{eq:g}
 g(s_1,s_2)=\sum_{n=1}^{\infty}\sum_{k=0}^n
 \frac{1}{n^2 2^n}\binom nk(-2)^k
 s_1^{,2n-2k}s_2^k.
\end{equation}
This series fails to converge absolutely at points of the symmetrized
bidisc $G_2:=\pi(\D^2)$.  Consequently it has no power-series convergence
domain containing $G_2$.
\end{theorem}

\begin{proof}
The monomials belonging to different $n$ in \eqref{eq:f} have different
total degrees.  Hence there is no cross-block coefficient cancellation and
\[
 \|f\|_1=\sum_{n=1}^{\infty}
 \frac{\|(z^2+w^2)^n\|_1}{n^2 2^n}
 =\sum_{n=1}^{\infty}\frac{2^n}{n^2 2^n}
 =\sum_{n=1}^{\infty}\frac1{n^2}<\infty.
\]
The function is symmetric.  The identity
$z^2+w^2=s_1^2-2s_2$ and the binomial theorem give \eqref{eq:g}.

Fix $r\in(1/\sqrt3,1)$ and take $(z,w)=(r,r)$.  Its image
$(s_1,s_2)=(2r,r^2)$ lies in $G_2$.  The sum of the absolute values of
all monomials in the $n$th block of \eqref{eq:g}, evaluated there, is
\begin{align*}
 &\frac1{n^2 2^n}\sum_{k=0}^n\binom nk
 2^k(2r)^{2n-2k}(r^2)^k\\
 &\hspace{3cm}=\frac{(4r^2+2r^2)^n}{n^2 2^n}
 =\frac{(3r^2)^n}{n^2}.
\end{align*}
These quantities do not even tend to zero because $3r^2>1$.
Also, a monomial with exponents $(2n-2k,k)$ has weighted degree
$(2n-2k)+2k=2n$, so distinct pairs $(n,k)$ give distinct monomials;
the calculation does not overcount coefficients which later combine.

A multivariable power series is absolutely and locally uniformly convergent
throughout the interior of its convergence set.  Thus \eqref{eq:g} cannot
converge on any domain containing the open set $G_2$.
\end{proof}

For the concrete choice $r=3/4$, the absolute $n$th block is
$(27/16)^n/n^2$.

\section*{What remains true: block factorization}

The negative answer concerns the formal \emph{monomial} series, not the
existence of a holomorphic factor.

\begin{proposition}[corrected positive statement]\label{prop:block}
For every $d<\infty$ and every
\[
 f(z)=\sum_{\alpha\in\N_0^d}c_\alpha z^\alpha
 \in W^+_{\rm sym}(\D^d),
\]
let $p_k=\sum_{|\alpha|=k}c_\alpha z^\alpha$ and let $Q_k$ be the
unique polynomial satisfying $p_k=Q_k\circ\pi$.  Then
\[
 F(s):=\sum_{k=0}^{\infty}Q_k(s)
\]
converges locally uniformly on $G_d:=\pi(\D^d)$, defines a holomorphic
function there, and satisfies $f=F\circ\pi$.
\end{proposition}

\begin{proof}
The finite symmetrization map $\pi:\D^d\to G_d$ is proper.  Given a compact
$K\subset G_d$, its preimage is compact in $\D^d$, so there is $\rho<1$
such that $|z_j|\le\rho$ for every $z\in\pi^{-1}(K)$ and every $j$.
For $s=\pi(z)\in K$,
\[
 |Q_k(s)|=|p_k(z)|
 \le\rho^k\sum_{|\alpha|=k}|c_\alpha|.
\]
The Weierstrass test proves uniform convergence on $K$.  Every partial sum
is a polynomial in $s$, so the limit is holomorphic.  Substitution gives
$F(\pi(z))=\sum_kp_k(z)=f(z)$.
\end{proof}

For Theorem \ref{thm:counter}, the block factor is
\[
 F(s_1,s_2)=\sum_{n=1}^{\infty}
 \frac{(s_1^2-2s_2)^n}{n^2 2^n}
 =\operatorname{Li}_2\!\left(\frac{s_1^2-2s_2}{2}\right).
\]
It is holomorphic on $G_2$ because
$(s_1^2-2s_2)/2=(z^2+w^2)/2$ has modulus less than one there.  At the
diagonal points used above, the polynomial blocks converge by cancellation,
while their monomial absolute values diverge.  This precisely separates
Proposition \ref{prop:block} from the source's formal-series hope.

\section*{Novelty and verification}

Exact-id, title, formal-series, elementary-symmetric-variable, and
symmetrized-polydisc searches were run in the local FA/Banach indexes and
primary arXiv sources through 2026-08-12.  No later answer or duplicate
packet was found.  The coefficient calculation was independently checked
with exact rational arithmetic; the source question and its example occur
on PDF pages 17--18.

\section*{References}

\begin{itemize}
\item A. Sasane, \emph{Banach algebras of symmetric functions on the
polydisc}, arXiv:2201.00183.
\end{itemize}

\end{document}
